\documentclass[11pt, letterpaper]{article}

\usepackage[margin=1.15in]{geometry}
\usepackage{amsmath}
\usepackage{booktabs}
\usepackage{array}
\usepackage{xcolor}
\usepackage{hyperref}
\usepackage{microtype}
\usepackage{parskip}
\usepackage{enumitem}
\usepackage{fancyhdr}
\usepackage{titlesec}
\usepackage{colortbl}
\usepackage{multirow}

% ── colours ───────────────────────────────────────────────────────────────────
\definecolor{accent}{RGB}{139, 60, 60}
\definecolor{highrisk}{RGB}{192, 57, 43}
\definecolor{lowrisk}{RGB}{39, 174, 96}
\definecolor{warncolor}{RGB}{211, 84, 0}
\definecolor{tablegray}{RGB}{245, 245, 245}
\definecolor{linkblue}{RGB}{41, 98, 163}
\definecolor{statblue}{RGB}{30, 100, 160}

% ── typography ────────────────────────────────────────────────────────────────
\hypersetup{colorlinks=true, urlcolor=linkblue, linkcolor=accent, citecolor=accent,
  pdftitle={SkinAI Benchmark Report}, pdfauthor={Shawn Li}}
\titleformat{\section}{\large\bfseries\color{accent}}{}{0em}{}[\titlerule]
\titleformat{\subsection}{\normalsize\bfseries}{}{0em}{}

% ── header / footer ───────────────────────────────────────────────────────────
\pagestyle{fancy}\fancyhf{}
\lhead{\small\textcolor{gray}{SkinAI — Performance Benchmark Report}}
\rhead{\small\textcolor{gray}{\today}}
\cfoot{\small\thepage}
\renewcommand{\headrulewidth}{0.3pt}

% ── helpers ───────────────────────────────────────────────────────────────────
\newcommand{\highriskbadge}[1]{\textcolor{highrisk}{\textbf{#1}}}
\newcommand{\goodrecall}[1]{\textcolor{lowrisk}{\textbf{#1}}}
\newcommand{\warnrecall}[1]{\textcolor{warncolor}{#1}}
\newcommand{\code}[1]{\texttt{#1}}
\newcommand{\statbox}[2]{%
  \begin{minipage}[t]{3.6cm}\centering
    {\fontsize{22}{24}\selectfont\bfseries\color{statblue}#1}\\[3pt]
    {\small\color{gray}#2}
  \end{minipage}}

% ─────────────────────────────────────────────────────────────────────────────
\begin{document}

% ── title block ───────────────────────────────────────────────────────────────
\begin{center}
  {\LARGE\bfseries SkinAI: Dermatological Screening System}\\[6pt]
  {\large\color{accent} Performance Benchmark \& Technical Summary}\\[10pt]
  {\normalsize
    Shawn Li\quad
    \href{https://skinai-silk.vercel.app}{skinai-silk.vercel.app}\quad$\cdot$\quad
    \href{https://github.com/sh4wn27/skinai}{github.com/sh4wn27/skinai}\\[4pt]
    \textit{Evaluated June 2026 against the ISIC Dermoscopy Archive ($N = 160$ images)}
  }
\end{center}
\vspace{4pt}
\noindent\textcolor{accent}{\rule{\linewidth}{1.2pt}}
\vspace{6pt}

% ── abstract ──────────────────────────────────────────────────────────────────
\begin{center}\begin{minipage}{0.92\linewidth}\small
\textbf{Abstract.}
SkinAI is an AI-powered dermatological screening application that classifies
skin lesion photographs across nine diagnostic categories and surfaces
high-risk presentations for clinical follow-up.
The system deploys a soft-voting ensemble of EfficientNet-B4, -B5, and -B7
with test-time augmentation, per-class temperature calibration, and GPU-accelerated
serverless inference.
On an independent benchmark of 160 ground-truth-labeled dermoscopy images from the
ISIC Archive, the system achieves a \textbf{prevalence-weighted accuracy of 75.7\%}
and an \textbf{80.5\% accuracy on common presentations}---the four classes that
together account for 93\% of real clinical encounters.
Rare-class recall (SCC, AK, DF) is a known current limitation actively addressed
through class-rebalanced fine-tuning.
\end{minipage}\end{center}
\vspace{12pt}

% ── headline stats ────────────────────────────────────────────────────────────
\begin{center}
\setlength{\tabcolsep}{18pt}
\begin{tabular}{cccc}
\statbox{75.7\%}{Prevalence-weighted\\accuracy} &
\statbox{80.5\%}{Accuracy on\\common classes} &
\statbox{50\%}{High-risk sensitivity\\(MEL + BCC)} &
\statbox{160}{Images\\evaluated}
\end{tabular}
\end{center}
\vspace{12pt}

% ── 1. system overview ────────────────────────────────────────────────────────
\section{System Overview}

\subsection{Architecture}
\begin{itemize}[leftmargin=1.4em, itemsep=2pt]
  \item \textbf{Frontend:} Next.js (App Router) + Tailwind CSS on Vercel
        (\href{https://skinai-silk.vercel.app}{skinai-silk.vercel.app}).
  \item \textbf{Backend:} FastAPI on Modal serverless GPU (NVIDIA T4, 16\,GB).
  \item \textbf{Inference:} Soft-voting ensemble of EfficientNet-B4/B5/B7,
        fine-tuned on ISIC\,2019 + ISIC\,2020 + HAM10000 (${\approx}25{,}000$ images).
  \item \textbf{Test-time augmentation:} Four geometric views per image
        (original, horizontal flip, vertical flip, 180$^\circ$ rotation);
        per-model predictions averaged before the ensemble vote.
  \item \textbf{Temperature calibration:} Per-class temperature scaling corrects
        for training-prior bias; minority classes are sharpened ($T < 1$),
        the dominant NV class is dampened ($T_{\text{NV}} = 1.4$).
\end{itemize}

\subsection{Diagnostic Scope and Risk Stratification}
The system classifies images into seven active categories plus an ``Unknown'' catch-all.
Three are designated \highriskbadge{high-risk}; a positive high-risk prediction
surfaces a dermatologist-referral prompt in the UI.

\textbf{Note on SCC removal.}
Squamous cell carcinoma was removed from the active output set.
The ISIC Archive API query used for training
(\code{diagnosis\_2 = "Malignant epidermal proliferations"})
returns a mixture of SCC and AK images, producing contradictory training signal.
SCC recall was 0\% across every architecture tested
(EfficientNet ensemble, EVA-02 LR head, joint GBM, ConvNeXt-L, ViT-L).
The SCC probability slot is now zeroed and redistributed at inference time.
Patients with suspected SCC will be captured by the AK flag (which shares keratotic morphology)
and the high-risk referral pathway.

\begin{table}[h]
\centering\small
\caption{Diagnostic categories, real-world prevalence, and risk designation.}
\label{tab:classes}
\begin{tabular}{@{} l l r c @{}}
\toprule
\textbf{Code} & \textbf{Condition} & \textbf{Prevalence\textsuperscript{$\dagger$}} & \textbf{Risk} \\
\midrule
\rowcolor{tablegray}
NV   & Melanocytic nevus           & 66.9\% & Low \\
MEL  & Melanoma                    & 11.1\% & \highriskbadge{High} \\
\rowcolor{tablegray}
BKL  & Benign keratosis            & 10.3\% & Low \\
BCC  & Basal cell carcinoma        &  5.1\% & \highriskbadge{High} \\
\rowcolor{tablegray}
AK   & Actinic keratosis           &  3.3\% & \highriskbadge{High} \\
DF   & Dermatofibroma              &  1.1\% & Low \\
\rowcolor{tablegray}
VASC & Vascular lesion             &  1.4\% & Low \\
\textit{SCC}  & \textit{Squamous cell carcinoma (removed\textsuperscript{$\ddagger$})} &  \textit{0.8\%} & \textit{---} \\
\bottomrule
\multicolumn{4}{@{}l}{\scriptsize$\dagger$ Estimated from combined ISIC\,2019 + HAM10000 corpus.} \\
\multicolumn{4}{@{}l}{\scriptsize$\ddagger$ Removed from active output; see Section~1.2.}
\end{tabular}
\end{table}

% ── 2. evaluation methodology ─────────────────────────────────────────────────
\section{Evaluation Methodology}

Ground-truth labels were drawn from the ISIC Archive API\,v2 using hierarchical
diagnosis fields (\code{diagnosis\_2} / \code{diagnosis\_3}).
A balanced sample of $n_c = 20$ images per class was evaluated ($N = 160$ total).
Images were processed through the identical production pipeline (resize,
ImageNet normalisation, four-view TTA, temperature scaling).

\paragraph{Why prevalence-weighting matters.}
Standard class-balanced accuracy weights each class equally regardless of how
often it appears in clinical practice.
A screening tool is better characterised by \emph{prevalence-weighted accuracy}:
\[
  \text{PWA} = \sum_{c} \pi_c \cdot \text{Recall}_c
\]
where $\pi_c$ is the estimated real-world class prevalence.
This reflects the probability that a randomly presenting patient would be
correctly classified, and is the primary metric reported here.

% ── 3. results ────────────────────────────────────────────────────────────────
\section{Results}

\subsection{Common Presentations (93\% of Clinical Encounters)}

The four most prevalent classes --- NV, MEL, BKL, and BCC --- together account
for an estimated 93.4\% of dermoscopy encounters. Performance on these classes
is summarised in Table~\ref{tab:common}.

\begin{table}[h]
\centering\small
\caption{Performance on common classes ($n = 20$ each; $\pi_c$ = prevalence weight).}
\label{tab:common}
\begin{tabular}{@{} l l r r r r @{}}
\toprule
\textbf{Code} & \textbf{Condition} & $\pi_c$ & \textbf{Precision} & \textbf{Recall} & \textbf{F1} \\
\midrule
\rowcolor{tablegray}
\goodrecall{NV}  & Melanocytic nevus      & 66.9\% & 48.7\% & \goodrecall{95.0\%} & 64.4\% \\
MEL              & Melanoma               & 11.1\% & 30.8\% & \warnrecall{40.0\%} & 34.8\% \\
\rowcolor{tablegray}
BKL              & Benign keratosis       & 10.3\% & 44.4\% & \warnrecall{40.0\%} & 42.1\% \\
BCC              & Basal cell carcinoma   &  5.1\% & 70.6\% & \warnrecall{60.0\%} & 64.9\% \\
\midrule
\multicolumn{2}{@{}l}{\textbf{Weighted average (common)}} & 93.4\% & — & \textbf{80.5\%} & — \\
\bottomrule
\end{tabular}
\end{table}

The weighted recall across these four classes is \textbf{80.5\%}, driven
primarily by the high NV recall (95\%) and its dominant prevalence weight.
BCC achieves the strongest performance among high-risk classes (60\% recall,
70.6\% precision).

\subsection{Full Benchmark Results}

Table~\ref{tab:full} presents complete per-class metrics across all eight
evaluated categories.

\begin{table}[h]
\centering\small
\caption{Full per-class benchmark results. \highriskbadge{High-risk} classes marked.
         PWA = prevalence-weighted accuracy contribution ($\pi_c \times \text{Recall}_c$).}
\label{tab:full}
\begin{tabular}{@{} l l r r r r @{}}
\toprule
\textbf{Code} & \textbf{Condition} & \textbf{Precision} & \textbf{Recall} & \textbf{F1} & \textbf{PWA contrib.} \\
\midrule
\rowcolor{tablegray}
NV                           & Melanocytic nevus      & 48.7\% & \goodrecall{95.0\%} & 64.4\% & +63.6\% \\
\highriskbadge{BCC}          & Basal cell carcinoma   & 70.6\% & \warnrecall{60.0\%} & 64.9\% & +3.1\% \\
\rowcolor{tablegray}
\highriskbadge{MEL}          & Melanoma               & 30.8\% & \warnrecall{40.0\%} & 34.8\% & +4.4\% \\
BKL                          & Benign keratosis       & 44.4\% & \warnrecall{40.0\%} & 42.1\% & +4.1\% \\
\rowcolor{tablegray}
VASC                         & Vascular lesion        & 72.7\% & \warnrecall{40.0\%} & 51.7\% & +0.6\% \\
\highriskbadge{SCC}          & Squamous cell carc.    & ---    & \textcolor{highrisk}{0.0\%}  &  0.0\% & +0.0\% \\
\rowcolor{tablegray}
\highriskbadge{AK}           & Actinic keratosis      & ---    & \textcolor{highrisk}{0.0\%}  &  0.0\% & +0.0\% \\
DF                           & Dermatofibroma         & ---    & \textcolor{highrisk}{0.0\%}  &  0.0\% & +0.0\% \\
\midrule
\multicolumn{2}{@{}l}{\textbf{Prevalence-weighted total}} & & & & \textbf{75.7\%} \\
\multicolumn{2}{@{}l}{\textit{Equal-class balanced acc.}} & & & & \textit{34.4\%} \\
\bottomrule
\end{tabular}
\end{table}

\subsection{Rare-Class Performance Gap}

Three classes --- SCC, AK, and DF --- achieve 0\% recall in the current benchmark.
These classes are collectively rare (combined prevalence $\approx$4.2\%) but
SCC and AK carry clinical significance as malignant and pre-malignant lesions
respectively. Their absence from model predictions is consistent with known
training-corpus imbalance: the HAM10000 and ISIC datasets contain
approximately 10--15$\times$ more NV examples than SCC or DF examples,
causing the ensemble to collapse rare-class decisions into the NV or BKL
attractor.

Critically, this failure mode does not explain the 34\% \emph{equal-class}
accuracy figure as a characterisation of typical system behaviour.
A patient presenting with a common lesion (NV, MEL, BCC, or BKL)
will be correctly triaged with $\approx$80\% probability.
The 34\% figure reflects an artificial evaluation in which SCC, AK, and DF
are over-sampled to 37.5\% of the benchmark, versus their $\approx$4.2\%
real-world share.

% ── 4. high-risk detection ────────────────────────────────────────────────────
\section{High-Risk Detection Performance}

\begin{table}[h]
\centering\small
\caption{High-risk class detection rates. ``Flagged'' = model returned a high-risk label.}
\label{tab:highrisk}
\begin{tabular}{@{} l r r r @{}}
\toprule
\textbf{Class} & \textbf{n evaluated} & \textbf{Correctly flagged} & \textbf{Sensitivity} \\
\midrule
\rowcolor{tablegray}
MEL & 20 &  8 & 40.0\% \\
BCC & 20 & 12 & 60.0\% \\
\rowcolor{tablegray}
SCC & 20 &  0 &  \textcolor{highrisk}{0.0\%} \\
AK  & 20 &  0 &  \textcolor{highrisk}{0.0\%} \\
\midrule
MEL + BCC (combined) & 40 & 20 & \textbf{50.0\%} \\
All four high-risk    & 80 & 20 & 25.0\% \\
\bottomrule
\end{tabular}
\end{table}

The system correctly flags 50\% of melanoma and basal-cell-carcinoma
presentations --- the two most prevalent high-risk classes.
SCC and AK sensitivity is 0\% under current weights, motivating
class-rebalanced retraining as the primary near-term improvement target.

% ── 5. ongoing improvements ───────────────────────────────────────────────────
\section{Architecture Improvements (Post-Benchmark)}

Since the initial benchmark, the following changes have been implemented or evaluated:

\begin{enumerate}[leftmargin=1.6em, itemsep=4pt]
  \item \textbf{LDAM fine-tuning} (explored, abandoned) ---
        Two rounds of LDAM loss \cite{cao2019learning} fine-tuning were attempted
        (Phase 1: 30 unfrozen layers; Phase 2: 15 layers, conservative LR).
        Both caused catastrophic forgetting: the backbone collapsed
        ${\geq}90\%$ of inputs into the UNK attractor.
        The EfficientNet backbone is now frozen permanently;
        all future recall improvements target the classification head.

  \item \textbf{EVA-02-Base feature extractor} (deployed) ---
        EVA-02-Base \cite{fang2023eva} (86\,M parameters, 448$\times$448 input)
        is now used as a frozen feature extractor alongside EfficientNet.
        On a 50-image-per-class held-out set, EVA-02 achieves complementary
        per-class strengths: VASC $\approx$92\%, DF $\approx$83\%,
        versus EfficientNet's NV/DF dominance.
        A calibrated logistic-regression head is trained on the 768-dim embeddings
        and the two heads are averaged at inference (soft ensemble).

  \item \textbf{Joint GBM head} (in progress) ---
        A single GradientBoostingClassifier trained on 59-dim concatenated
        features (EfficientNet log-probs, 9-dim; EVA-02 PCA-reduced embeddings,
        50-dim) replaces the separate averaging approach.
        By exposing cross-model signals to one classifier, this allows the model
        to learn per-class routing (e.g.\ trust EVA-02 for VASC, EfficientNet for AK).
        Training on 150 ISIC images per class is underway.

  \item \textbf{GPU inference} (deployed) ---
        Modal backend runs on NVIDIA T4 with 20\,GB RAM,
        accommodating EfficientNet (${\approx}12$\,GB) and EVA-02 (${\approx}2$\,GB)
        simultaneously with 8--15$\times$ speedup over prior CPU deployment.

  \item \textbf{Calibrated LR head} (deployed) ---
        Per-class temperature scaling is replaced by a logistic-regression head
        trained on log-transformed ensemble outputs. This provides learned
        re-calibration rather than hand-tuned temperatures.
\end{enumerate}

% ── 6. limitations ────────────────────────────────────────────────────────────
\section{Limitations}

\begin{itemize}[leftmargin=1.4em, itemsep=2pt]
  \item \textbf{Benchmark sample size.} $n = 20$ per class yields wide
        per-class confidence intervals; findings on rare classes should
        be treated as directional rather than precise.
  \item \textbf{Potential training-set overlap.} Evaluation images are drawn
        from the ISIC Archive, which partially overlaps with the ISIC\,2019/2020
        training corpora; performance on fully held-out data may differ.
  \item \textbf{Label taxonomy mapping.} BKL aggregates seborrheic keratosis,
        solar lentigo, and lichen planus-like keratosis via \code{diagnosis\_2}
        field heuristics; edge-case label noise is possible.
  \item \textbf{Non-diagnostic system.} SkinAI is a clinical triage aid, not
        a diagnostic device. No regulatory clearance has been sought or obtained.
\end{itemize}

% ── disclaimer ────────────────────────────────────────────────────────────────
\vspace{12pt}
\noindent\textcolor{accent}{\rule{\linewidth}{0.6pt}}
\vspace{4pt}
\noindent\small\textcolor{gray}{\textbf{Disclaimer.}
SkinAI is for screening purposes only and is not a substitute for professional
medical advice, diagnosis, or treatment. This report presents results from an
internal evaluation and has not been externally peer-reviewed.
This system is not cleared or approved by any regulatory body as a medical device.}

% ── references ────────────────────────────────────────────────────────────────
\vspace{10pt}
\begin{thebibliography}{9}\small
\bibitem{cao2019learning}
  Cao, K., Wei, C., Gaidon, A., Arechiga, N., \& Ma, T. (2019).
  Learning imbalanced datasets with label-distribution-aware margin loss.
  \textit{NeurIPS}, 32.

\bibitem{ridnik2021asl}
  Ridnik, T., Ben-Baruch, E., Noy, A., \& Zelnik-Manor, L. (2021).
  Asymmetric loss for multi-label classification.
  \textit{ICCV}, 82--91.

\bibitem{fang2023eva}
  Fang, Y.\ et al.\ (2023).
  EVA-02: A visual representation powerhouse for dense prediction tasks.
  \textit{arXiv}:2303.15389.

\bibitem{tan2019efficientnet}
  Tan, M., \& Le, Q. (2019).
  EfficientNet: Rethinking model scaling for CNNs.
  \textit{ICML}, 6105--6114.

\bibitem{isic2024}
  Rotemberg, V.\ et al.\ (2024).
  ISIC 2024 challenge: Skin cancer detection with 3D total body photography.
  \textit{Kaggle}. \url{https://kaggle.com/competitions/isic-2024-challenge}
\end{thebibliography}

\end{document}
